\documentclass[../../main.tex]{subfiles}
\begin{document}

$\vec{A}$はその大きさが$a$であるような定ベクトルであり，
また，各$k$(=1，2，3，4，5，6)について，
次のような成分をもつベクトル$\vec{B}_k$が与えられている．
$\displaystyle \vec{B}_k = \left( \cos\dfrac{k\pi}{3},\sin\dfrac{k\pi}{3} \right)$．

一つのサイコロを投げて，$k$の目が出たとき，
ベクトル$\vec{X}=\vec{A}+\vec{B}_k$をとるものとする．
この場合，ベクトル$\vec{X}$の大きさの平方の期待値を，$a$を用いて表わせ．

\end{document}
